\subsection*{Simulation implementation}
\label{sec:methods:sim}

All simulations were implemented in Python 3.x using NumPy~1.x. Random seed was fixed at 2026
for reproducibility. $N = 500$ agents were simulated per condition. Source code is available at
\url{https://github.com/example/bee-precision-model} (placeholder).

\subsection*{Precision distribution parameterization}
\label{sec:methods:params}

Precision distributions were derived from published bee reversal learning $d'$ data
\cite{AvarWeber2012}. The reported $d'$ range is approximately $[0.3, 2.1]$ for individual
honeybees, normalized to $[0.2, 0.95]$ by the transformation
$\precision = 0.2 + (d'/2.5) \times 0.75$.
The resulting empirical distribution was approximated by:
\begin{itemize}
  \item Normal condition: Beta$(5, 2) \times 0.75 + 0.20$ (mean $\approx 0.64$)
  \item Disrupted condition: Beta$(2, 4) \times 0.60 + 0.10$ (mean $\approx 0.30$)
  \item Nurse/forager conditions: as normal but with adjusted caste-identity priors
\end{itemize}

\subsection*{Analytical derivation}
\label{sec:methods:deriv}

The precision optimum was computed numerically on a grid of 999 equally spaced points
$\precision \in [0.01, 0.99]$. The analytical result (inverted-U, interior maximum) was
verified symbolically by checking the sign of $d\metaacc/d\precision$ at the boundaries.

\subsection*{GCA factor extraction}
\label{sec:methods:gca}

The GCA factor was defined as the fraction of variance explained by the leading eigenvalue of the
$3 \times 3$ covariance matrix of [tool use, caste learning, GCA score]. Eigendecomposition was
performed via \texttt{numpy.linalg.eigvalsh}.

\subsection*{Statistical analysis}
\label{sec:methods:stats}

All correlations are Pearson's $r$ computed over $N = 500$ individual agents per condition.
No correction for multiple comparisons was applied, as all reported correlations are directional
predictions from the model rather than exploratory comparisons.
